\section{迭代检索与多步推理}
\label{sec:iterative}

\subsection{检索如何随中间结果展开}

多跳问题的后续查询往往依赖前一步答案。询问“某部获奖电影导演的生日”时，系统只有先确定电影和导演，才能有效检索出生日期。增加第一次检索的文档数量未必能弥补这一依赖，因为原始问题尚未提供第二次查询所需的实体。

IRCoT 将思维链步骤与检索交替执行，让已有推理形成下一步查询，再用新证据支持后续推理~\citep{trivedi2023ircot}。Iter-RetGen 利用模型输出揭示尚缺的信息，将已有生成与原始问题结合用于下一轮检索，随后根据新证据改写输出~\citep{shao2023iterretgen}。两者都保留了一个不断更新的中间状态，区别主要在于用显式推理步骤还是已有生成结果指导检索。

ReAct 将推理轨迹与外部行动交错，使模型能够根据观察结果调整计划或处理异常~\citep{yao2023react}。在这一框架下，搜索成为一种工具调用，检索可以与其他操作共同组成任务流程。Agentic RAG 延续这一方向，引入规划、反思、工具使用和多智能体协作~\citep{singh2025agentic}。行动空间由此扩大，但证据评价、预算控制和停止条件仍需要具体机制来实现。

\subsection{多轮过程中的错误传播与停止}

迭代为系统提供了修正机会，也让中间错误有了更多传播路径。第一轮若识别错实体，后续查询可能持续收集与该实体相关却无助于原问题的文档；旧输出一直保留在上下文中时，早期无依据陈述还可能被后续步骤当作前提。重复文档也会带来另一个问题：数量上的增加容易被误当成多个来源的支持。因此，每一轮检索需要回答“新增了什么证据”，而不只是产生了多少新文本。

停止条件直接影响错误和成本。固定最大轮数便于实现，却无法区分证据已经齐全与检索陷入重复两种情况。更有信息量的依据包括关键事实是否得到支持、上一轮是否增加了新来源、冲突是否得到处理，以及剩余预算是否允许继续。IRCoT 和 Iter-RetGen 展示了迭代检索的收益，部署时还需要将这些过程信号与最大轮数、词元上限和超时约束结合~\citep{trivedi2023ircot,shao2023iterretgen}。

流程复杂度应与任务需求匹配。单跳事实问答通常无需规划器或多智能体协作；跨源比较、长期搜索和动态计划更能发挥 Agentic 方法的作用。实际系统可以保留几条不同长度的路径，让简单请求直接结束，证据失败时进入纠正流程，只有复杂依赖才启用多轮推理。
